\section{Tính năng và Chức năng của ứng dụng}
\label{sec:application_features}

\subsection{Các tính năng cốt lõi}
Sản phẩm khả thi tối thiểu (MVP) tập trung vào việc cung cấp luồng người dùng thiết yếu: khám phá các món ăn, lưu các sở thích, xem lại các lựa chọn và điều hướng đến điểm đến.
\begin{itemize}
    \item \textbf{Vuốt để khám phá (Cơ chế Swipe-to-Match):} Đây là cơ chế khám phá chính. Hệ thống lấy các món ăn gần đó dựa trên vị trí của người dùng và hiển thị chúng trong Giao diện Chế độ xem Đơn (Single-View Interface). Người dùng vuốt sang trái để bỏ qua hoặc vuốt sang phải để lưu món ăn. Cơ chế này thúc đẩy việc ra quyết định nhanh chóng và trực quan.
    \item \textbf{Quản lý danh sách "Muốn thử" (Bộ sưu tập cá nhân):} Một không gian dành riêng nơi người dùng có thể xem lại và quản lý các món ăn họ đã lưu trước đó thông qua việc vuốt. Nó cho phép người dùng so sánh các lựa chọn, cân nhắc lại các lựa chọn và xóa các mục, chuyển đổi giai đoạn khám phá thành một quy trình ra quyết định có thể hành động.
    \item \textbf{Thẻ thông tin nhà hàng (Restaurant Info Card):} Khi tương tác với một món ăn, ứng dụng sẽ hiển thị một thẻ toàn diện bao gồm hình ảnh thực phẩm chất lượng cao, tóm tắt thực đơn, giờ mở cửa và các đánh giá nổi bật. Điều này cung cấp một cái nhìn tổng quan nhanh chóng mà không yêu cầu chuyển đổi nhiều màn hình.
\end{itemize}

\subsection{Các chức năng nâng cao}
Để nâng cao trải nghiệm người dùng ngoài việc khám phá cơ bản, một số chức năng nâng cao đã được tích hợp vào nền tảng:
\begin{itemize}
    \item \textbf{Bộ lọc thông minh (Smart Filters):} Người dùng có thể thu hẹp kết quả tìm kiếm một cách chính xác dựa trên nhu cầu thời gian thực bằng cách lọc theo bán kính (khoảng cách), mức giá, loại ẩm thực và xếp hạng.
    \item \textbf{Tích hợp Bản đồ \& Điều hướng:} Người dùng có thể xem các nhà hàng đã thích trên bản đồ hoặc điều hướng trực tiếp đến vị trí bằng một lần chạm thông qua việc tích hợp Google Maps hoặc Apple Maps, giúp tối ưu hóa lộ trình di chuyển.
    \item \textbf{Tích hợp API bên thứ ba:} Để tăng uy tín dữ liệu và cung cấp cái nhìn khách quan, ứng dụng tích hợp với các nguồn bên ngoài (ví dụ: Google Places, Yelp) để hiển thị điểm số và đánh giá đã được xác minh.
\end{itemize}

\subsection{Hướng phát triển trong tương lai}
Trong các giai đoạn phát triển tiếp theo, ứng dụng sẽ bổ sung các tính năng nâng cao nhằm tăng tính kết nối cộng đồng:
\begin{itemize}
    \item \textbf{Khớp nhóm (Group Match):} Cơ chế khớp chỉ xảy ra khi mọi người trong nhóm cùng vuốt sang phải, giúp giải quyết các xung đột lựa chọn trong nhóm một cách công bằng.
    \item \textbf{Gamification (Vòng quay may mắn):} Biến việc chọn nhà hàng thành một trò chơi xã hội thú vị, hỗ trợ giải quyết xung đột khi nhóm không thể quyết định chung.
    \item \textbf{Tour ẩm thực (Food Tours):} Ứng dụng hỗ trợ các hành trình được tuyển chọn để biến một bữa tối tiêu chuẩn thành một cuộc phiêu lưu nhỏ, kích thích sự tò mò và tăng thời gian người dùng tương tác với nền tảng.
\end{itemize}
